\section{Cold estimates up to the retirement times}
\label{sec:cold}

A jet remains a separate color until \(L_s=s^{-1/2}\), which tends to
infinity as its scale tends to zero.  We now control it throughout that
interval.  The essential facts are that the free jet remains in a
narrow spatial tube, its self-collisions create only a small Gaussian
remainder, and the active tubes stay disjoint on the torus.  These are
long-time versions of the cold estimates in
\cite[Section 4]{OpenAIPeriodic}.  We give the details because both the
long deadlines and the change of lattice matter for global continuation.
All constants may depend on the four fixed parameters
\eqref{con:parameters}, on \(J\), and on the fixed choice of \(d_0\),
but not on \(N,b,T\) or the size of the hot component.

\subsection{Disjoint spatial enlargements}

For \(i=(j,\sigma,n)\), put
\begin{equation}
 \widehat E_j^\sigma
   =\{y:\dist(y,E_j^\sigma)\le(c_\eta/4)s_j^2\},\qquad
 \mathcal O_i(t)=n+\widehat E_j^\sigma\times
     [-a-2t-b_{\mathrm{sp}},a+2t+b_{\mathrm{sp}}].
 \label{cold:enlargements}
\end{equation}
These sets enlarge the spatial supports for estimates; the original
stripe masks in \(c_i^0\) are always retained when taking a limit.

\begin{lemma}[Disjointness before retirement]\label{cold:geometry}
For sufficiently small \(s_0\), the sets \(\mathcal O_i(t)\) for
active labels are pairwise disjoint in the periodic lift.  More
precisely, \(\mathcal O_i(t)\cap\mathcal O_{i'}(t)=\varnothing\)
for distinct labels whenever
\(t\le\min\{L_{s_i},L_{s_{i'}}\}\).  Also
\[
 \widehat E_j^\sigma\subset\{s_j<|y|<e^\eta s_j\}.
\]
\end{lemma}
\begin{proof}
Within a fixed copy, the radial margins are larger than the
corresponding enlargement radii.  Adjacent radial cores have gap
\(c_\eta(s_j^2+s_{j+1}^2)\), whereas the sum of their enlargement
radii is one quarter of that gap.  Opposite stripe patterns at a
single scale have gap \(s_j^2/4\), larger than the sum
\(c_\eta s_j^2/2\) of their enlargement radii.  This proves the
assertions for one copy.

For distinct copies let \(S_* =\max\{s_i,s_{i'}\}\).  An intersection
would force
\[
 |\bar n-\bar n'|\le4S_*,\qquad
 |n-n'|\le C(1+t)\le C(1+S_*^{-1/2}).
\]
By \eqref{con:diophantine}, the first quantity is at least
\(c/(1+|n-n'|)\ge c' S_*^{1/2}\).  For sufficiently small \(s_0\)
this contradicts \(4S_*\).  The choice is uniform over all pairs of
scales and all lattice copies.
\end{proof}

At a point \((t,x)\), the \emph{main label} is the unique active
label whose enlargement contains \(x\), if one exists.  Every other
active label is called nonmain at that point.  Write \(\mathcal O(t)\)
for the union, in the torus, of the active enlargements.  Including
additional scales absent from the finite approximation is harmless,
provided their deadlines have not passed.  At a retirement endpoint,
these definitions may use the colors just before they are transferred.

\subsection{A supersolution for one color}

For all \(s>0\), signs \(\sigma\), and \(t\ge0\), define the
axial profile \(U_{s,\sigma}(t,z)\) by
\begin{equation}
 U_{s,\sigma}(t,z)
       =a_s^\sigma\bigl(z-tU_{s,\sigma}(t,z)\bigr).
 \label{cold:profile}
\end{equation}
This is uniquely and smoothly defined: the map
\(\tau\mapsto\tau+t a_s^\sigma(\tau)\) is onto and has derivative
at least one.  In particular \(|U_{s,\sigma}|\le1+\alpha\).
For a fixed label, abbreviate \(s=s_j\), \(L=L_s\),
\(U=U_{s,\sigma}(t,z_c)\), and set
\[
 \tau=z_c-tU,\qquad \xi=v-Ue_3,\qquad
 D_y=\dist(y_c,E_j^\sigma),\qquad
 D_z=\dist(\tau,[-a,a]).
\]
The free transport and its majorant are
\begin{equation}
 f_i(t,x,v)=c_i^0(x-tv,v),
 \label{cold:free-core}
\end{equation}
\begin{equation}
 m_i(t,x,v)=C_m\frac K{s}d_s^{-3}e^{t/L}
 \exp\left[-\frac{(2-t/L)|\xi|^2}{d_s^2}
           -\frac{D_y^2+D_z^2}{8d_s^2Lt}\right],\qquad 0<t\le L.
 \label{cold:barrier-definition}
\end{equation}
Its value at time zero is the trace along free characteristics.

\begin{lemma}[Barrier estimates]\label{cold:barrier}
A fixed sufficiently large \(C_m\) gives, for \(0\le t\le L_s\),
\begin{align}
 f_i&\le m_i,\qquad
 \D m_i\ge \frac1{L_s}
       \left(1+\frac78\frac{|\xi|^2}{d_s^2}\right)m_i,
 \label{cold:barrier-derivative}\\
 P(m_i,m_i)&=m_i\nu_{m_i}
       \le C\frac K{s}(d_s+|\xi|)m_i.
 \label{cold:barrier-gain}
\end{align}
The differential inequality holds almost everywhere along flights.
For sufficiently small \(s_0\),
\(f_i(t,x,v)>0\) implies \(x\in\mathcal O_i(t)\) throughout the
same time interval.
\end{lemma}
\begin{proof}
Put \(A=(a_s^\sigma)'(\tau)\ge0\).  Differentiation of
\eqref{cold:profile} gives
\[
 U_z=\frac A{1+tA}\ge0,\qquad U_t=-UU_z,\qquad
 \D U=(v_3-U)U_z,\qquad \D\tau=\frac{v_3-U}{1+tA}.
\]
Thus \(\D|\xi|^2=-2U_z\xi_3^2\).  The distance functions are locally
Lipschitz along every flight, and
\[
 |\D D_y|\le|\bar\xi|,\qquad |\D D_z|\le|\xi_3|.
\]
With \(r=(D_y,D_z)\), differentiation of the logarithm of the barrier
therefore yields
\[
 \D\log m_i
 =\frac1L+\frac{|\xi|^2}{Ld_s^2}
   +\frac{2(2-t/L)U_z\xi_3^2}{d_s^2}
   +\frac{|r|^2}{8d_s^2Lt^2}
   -\frac{r\cdot\D r}{4d_s^2Lt}.
\]
The last two terms are at least
\(-|\D r|^2/(8d_s^2L)\), by completing a square.  Since
\(|\D r|\le|\xi|\) and \(2-t/L\ge1\), this proves the derivative
estimate.  The large axial derivative \(U_z\) has a favorable sign.

For a flight starting in the core at \(x^0\), write
\(v=a_s^\sigma(z_c^0)e_3+d_sw\), \(|w|\le1\).
At \(x=x^0+tv\), the increasing function
\(p\mapsto p-a_s^\sigma(z_c-tp)\) has derivative at least one,
value \(d_sw_3\) at \(p=v_3\), and value zero at \(p=U\).
Consequently
\[
 |\xi|\le d_s,\qquad
 D_y\le td_s|\bar w|,\qquad D_z\le td_s|w_3|.
\]
The distance exponent in \eqref{cold:barrier-definition} is at most
\(t/(8L)\), and its velocity exponent is at most two.  A fixed
\(C_m\), depending only on \(\|J\|_\infty\), thus gives \(f_i\le m_i\).
Since \(L_sd_s=d_0s^{19/2}\ll s^2\), the transverse displacement is
less than the enlargement radius.  The axial speed is below two, so
the flight also stays inside the axial interval of
\eqref{cold:enlargements}.

At fixed \((t,x)\), the majorant is a scalar multiple of a Maxwellian
centered at \(Ue_3\).  Detailed balance gives
\(P(m_i,m_i)=m_i\nu_{m_i}\).  Its mass is at most \(CK/s\), and its
first moment about \(Ue_3\) is at most \(CKd_s/s\), uniformly for
\(t\le L_s\).  The inequality
\(|v-p|\le|\xi|+|p-Ue_3|\) gives \eqref{cold:barrier-gain}.

Finally, along a flight starting on the closed initial spatial
support, both distance functions vanish initially and are locally
Lipschitz, so their squared-distance-over-time terms tend to zero.
Away from that support the barrier tends to zero.  These observations
give its path trace and allow the differential inequalities to be
integrated from zero.  Boundary conventions for the masks change only
null sets; no spatial indicator is differentiated.
\end{proof}

\begin{proposition}[Control of every active color]\label{cold:control}
There is a fixed constant \(C_*\) such that, with
\[
 \epsilon_s=C_*L_sd_s/s,\qquad
 \mathcal U_i=f_i+\epsilon_s m_i,
\]
and sufficiently small \(s_0\), every finite colored evolution obeys
\begin{equation}
 0\le c_i(t,x,v)\le\mathcal U_i(t,x,v)\qquad (0\le t\le L_{s_i})
 \label{cold:upper}
\end{equation}
as long as that color is active and the solution exists.
If, on an interval beginning at \(t_1\), one also has
\(\nu_F^b\ge\kappa>0\), then
\begin{equation}
 c_i(t,x,v)\le e^{-\kappa(t-t_1)}\mathcal U_i(t,x,v)\qquad
                   (t_1\le t\le L_{s_i})
 \label{cold:decay}
\end{equation}
throughout that interval.  None of the constants in
\eqref{cold:upper} uses a bound on the bath or hot density.
\end{proposition}
\begin{proof}
Reduce \(s_0\) so that \(\epsilon_s\le1\) at every scale.  Then
\(\mathcal U_i\le2m_i\), and Lemma~\ref{cold:barrier} gives
\[
 P^b(\mathcal U_i,\mathcal U_i)
 \le4C\frac K{s}(d_s+|\xi|)m_i,\qquad
 \D\mathcal U_i\ge\frac{\epsilon_s}{L_s}
          \left(1+\frac78\frac{|\xi|^2}{d_s^2}\right)m_i.
\]
Since \(1+z^2\ge c(1+z)\), a fixed large \(C_*\) makes
\(\D\mathcal U_i\ge P^b(\mathcal U_i,\mathcal U_i)\).
Dropping the nonnegative loss in the equation for \(c_i\) permits
positive gain comparison from the initial trace.

For completeness, this comparison is valid in the measurable solution
class of Proposition~\ref{con:finite-evolution}.  At a fixed
approximation \(m_i\le C_sH\) throughout the interval, and
\(z_i=(c_i-\mathcal U_i)_+\) has a pathwise upper integral comparison
from zero with source
\(2P^b(\mathcal U_i,z_i)+P^b(z_i,z_i)\).
The bilinear \(H\)-norm estimate and Gronwall's inequality give
\(\|z_i\|_H=0\).  One may first integrate from a positive time and
then use the trace established in Lemma~\ref{cold:barrier}.

For the second assertion,
\(e^{-\kappa(t-t_1)}\mathcal U_i\) is a supersolution for
\(\D+\kappa\): its linear derivative has factor
\(e^{-\kappa(t-t_1)}\), whereas its quadratic gain has the smaller
factor \(e^{-2\kappa(t-t_1)}\).  Its value at \(t_1\) dominates
\(c_i(t_1)\).  The same comparison proves \eqref{cold:decay}.
\end{proof}

\subsection{Velocity measures and sums over colors}

The preceding comparison bounds a narrow free core and a small
remainder separately.  The next lemma identifies the mass of that
core at a fixed spatial point.  Axial expansion produces the factor
\((t+s)^{-1}\), even though the initial velocity-integrated density is of order \(s^{-1}\).

\begin{lemma}[Core formula and sheet amplitudes]\label{cold:core-amplitude}
At a fixed \((t,x)\), the velocity measure \(f_i(t,x,p)\dd p\) is the
pushforward of
\begin{equation}
 \frac{K S(\gamma(z_c-tp_3)/s)}
          {s+t\alpha\gamma S(\gamma(z_c-tp_3)/s)}\,J(w)\dd w,\qquad
 \bar p=d_s\bar w,\quad
 p_3=a_s^\sigma(z_c-tp_3)+d_sw_3,
 \label{cold:core-formula}
\end{equation}
with the initial spatial masks evaluated at \(x-tp\).
For \(t\le L_s\), it is dominated by probability sheets with total
amplitude \(C/(t+s)\).  The remainder
\(\epsilon_s m_i\dd p\) has sheet amplitude at most
\begin{equation}
 C L_sd_s/s^2\le Cs^{15/2}\le C/(t+s).
 \label{cold:remainder-amplitude}
\end{equation}
Thus the full upper bound \(\mathcal U_i\dd p\) has sheet amplitude \(C/(t+s)\).
In particular, for every fixed \(k\ge0\), a main color satisfies
\begin{equation}
 \int\langle p\rangle^k c_i(t,x,p)\dd p\le\frac{C_k}{t+s},\qquad
 \int_{|p|>R}\langle p\rangle^k c_i(t,x,p)\dd p
       \le\frac{C_k e^{-c_kR^2}}{t+s}\quad(R\ge3).
 \label{cold:moments}
\end{equation}
\end{lemma}
\begin{proof}
The transverse Jacobian is \(d_s^2\), while differentiation of the
implicit equation for \(p_3\) gives
\[
 \frac{\partial p_3}{\partial w_3}
 =\frac{d_s}{1+t(a_s^\sigma)'(z_c-tp_3)}.
\]
Multiplying by the density in \eqref{con:jet-datum} proves
\eqref{cold:core-formula}.  Its coefficient is at most
\(C/(s+t)\), since \(0\le S\le1\) and \(\alpha\gamma>0\).
A fixed product Gaussian dominates \(J\).  Moreover, the monotonicity
argument in Lemma~\ref{cold:barrier} gives
\(|p_3-U|\le d_s|w_3|\).  The Gaussian parameter estimate
(Lemma~\ref{ker:parameter-sheets}) therefore supplies the asserted
sheet domination, with constants independent of the axial compression.

The remainder is a Gaussian centered at \(Ue_3\), of width comparable
to \(d_s\), whose mass is at most
\(C\epsilon_s K/s=C L_sd_s/s^2\).  Its transverse Gaussian width is
within the sheet class and its axial exponential moment is uniformly
bounded because \(|U|\le1+\alpha\).  Since
\(L_sd_s/s^2=d_0s^{15/2}\) and \(t\le s^{-1/2}\), this also gives the
last inequality of \eqref{cold:remainder-amplitude}.  Uniform Gaussian
moments and tails, together with \eqref{cold:upper}, prove
\eqref{cold:moments}.
\end{proof}

\begin{lemma}[Nonmain tails]\label{cold:nonmain}
There is a function \(\theta_0(s_0)\to0\) as \(s_0\downarrow0\)
such that, uniformly at all times of existence,
\begin{equation}
 \sum_{\substack{i\ \mathrm{active}\\x\notin\mathcal O_i(t)}}
       \mathcal U_i(t,x,v)
       \le\theta_0(s_0)e^{-4|v|^2}.
 \label{cold:nonmain-bound}
\end{equation}
For \(0<t\le1\), and any fixed \(A,B>0\), the coefficient can instead
be taken to be \(C_{A,B}t^A s_0^B\).  Both assertions remain valid
for an arbitrary subset of the active labels.
\end{lemma}
\begin{proof}
Free cores vanish off their enlargements.  For a nonmain label either
\(D_y\ge(c_\eta/4)s^2\), or
\(|z_c|>a+2t+b_{\mathrm{sp}}\), in which case
\(D_z\ge b_{\mathrm{sp}}\) because \(|U|<2\).
Completing the velocity square in the remainder gives
\[
 \mathcal U_i(t,x,v)
 \le C s^{-31}e^{-4|v|^2}
       \exp\left(-\frac{D_y^2+D_z^2}{8d_s^2L_st}\right).
\]
Here the Gaussian rate is at least eight once the widths are fixed
small, and its centers stay bounded.  Half the distance exponential
is at most
\(\exp[-c/(s^{16}L_st)]\).  Since \(t\le L_s\), this is bounded by
\(e^{-c/s^{15}}\), which absorbs every polynomial in \(1/s\).
If also \(t\le1\), it absorbs arbitrary powers of \(t\) and \(s\).

It remains to sum the other half over lattice copies.  Uniformly,
\[
 D_y\ge(|y_c|-C)_+,\qquad
 D_z\ge(|z_c|-C(1+t))_+.
\]
The Gaussian distance scale \(d_s\sqrt{L_st}\) is bounded, and the
central spatial region has diameter \(O(1+L_s)\).  The lattice has
unit separation, hence at most \(C(R+1)^3\) points in a ball of
radius \(R\).  Summing in successive shells gives at worst a fixed
polynomial loss in \(1/s\).  The extracted exponential absorbs this
loss and the factor \(s^{-31}\).  Finally sum over the geometric
scales.  This proves both claims, including their stated uniformity.
\end{proof}

\begin{lemma}[Size at retirement and volume of active tubes]
\label{cold:retirement-size}
For some fixed constants \(C,C'>0\),
\begin{equation}
 \left\|\sum_{i:\,s_i=s}\mathcal U_i(L_s)\right\|_H
       \le C s^{-C'}.
 \label{cold:retirement-bound}
\end{equation}
Moreover, uniformly over the active label sets and all times,
\begin{equation}
 |\mathcal O(t)|\le\sum_{j\ge1}C s_j^2(1+L_{s_j})
                    \le C s_0^{3/2}.
 \label{cold:tube-volume}
\end{equation}
\end{lemma}
\begin{proof}
For \eqref{cold:retirement-bound}, use \(f_i\le m_i\) and the same
Gaussian lattice sum as in Lemma~\ref{cold:nonmain}, without extracting
a small distance factor.  The prefactor and the number of central
copies are polynomial in \(1/s\).  Completing the velocity square
bounds the remaining velocity factor by \(Ce^{-4|v|^2}\le CH(v)\).
For \eqref{cold:tube-volume}, the transverse area at scale \(s_j\)
is \(O(s_j^2)\), and its axial length while active is
\(O(1+L_{s_j})\).  Projection onto the torus does not increase volume.
Summation gives \(C\sum_j(s_j^2+s_j^{3/2})\le Cs_0^{3/2}\).
\end{proof}

All pointwise assertions above mean essential bounds.  They also hold
at the time sections defined by the finite path equations: free
coordinates preserve measure, and Fubini's theorem supplies a
full-measure set of paths on which the inequalities can be integrated.
For each fixed terminal time this gives a full-measure endpoint set.
Countably many approximation indices, lattice copies, and interval
choices can be handled simultaneously.  This convention will also be
used for the path estimates in the next section.
